% █▀█ █▀█ █▀▀ ▄▀█ █▀▄▀█ █▄▄ █   █▀▀
% █▀▀ █▀▄ ██▄ █▀█ █ ▀ █ █▄█ █▄▄ ██▄

% NOTE: The packages (and related commands) in this preamble are loaded following the bellow ordering.
% Packages related to:
	% 1. Document canvas (i.e.: page geometry);
	% 2. Typeface system (i.e.: used fonts);
	% 3. Document structure (i.e.: headers formatting);
	% 4. Floats (i.e.: non-textual elements);
	% 5. Micro-Layout (i.e.: "inline" adjustments);
	% 6. Finalization (i.e.: packages best loaded last).

%::::::::::::::::::::::::::::::::::::::::::::::::::::::::::::::::::::::::::::::::
%                                   1. CANVAS                                    
%           Define document's physical space plus some basic utilities           
%................................................................................
% CANVAS
\usepackage[
	a4paper,
	margin = 2cm,
]{geometry}			% Set page geometry
\usepackage[brazil]{babel}	% Adjust document to the language it's written in

% BASIC UTILITIES
\usepackage{csquotes}		% Language-sensitive quoting;
				% 'csquotes' works in tandem with 'babel'
\usepackage[dvipsnames]{xcolor}	% Access to 68 named colors
\usepackage{moresize}		% Additional font sizes
\usepackage[normalem]{ulem}	% More underline styles (and aligns bottoms of multiple underlines)
				% 'normalem' prevents underlines from altering \emph


%::::::::::::::::::::::::::::::::::::::::::::::::::::::::::::::::::::::::::::::::
%                               2. TYPEFACE SYSTEM                               
%    Define all fonts - for text, algorithms and math - used in the document     
%................................................................................
% 2.1 SYSTEM BACKEND
\usepackage{fontspec}		% Native Unicode-to-OpenType mapping;
				% Replaces legacy TeX font-enconding hacks;
				% DOCUMENT MUST NOW BE COMPILED WITH Lua OR XeLaTex

\usepackage{amsmath}		% AMS's package for mathematical typesetting
\usepackage{mathtools}		% Extends amsmath's functionalities (i.e.: a superset)
\usepackage[
	math-style = ISO,
	bold-style = ISO
]{unicode-math}			% Must be loaded after 'mathtools';
				% DOCUMENT MUST NOW BE COMPILED WITH Lua OR XeLaTex

%. . . . . . . . . . . . . . . . . . . . . . . . . . . . . . . . . . . . . . . . . 
% 2.2 SYSTEM DEFINITION
% 2.2.1 TEXT FONTS
\setmainfont{Tex Gyre Pagella}[Scale = 1.0]
\setsansfont{URW Classico}[
	Scale = MatchLowercase,	% Ensure main and sans fonts are at the same scale
]
\setmonofont{Inconsolata}[Scale=MatchLowercase]

% 2.2.2 MATH FONTS
\setmathfont{Libertinus Math}[Scale=MatchLowercase]
% Manually pick distinct calligraphic (Chancery hand) and script (Round hand) fonts for a greater variety of symbols
\setmathfont{XITS Math}[
	Scale 		= MatchLowercase,
	range		= {cal,bfcal},
	StylisticSet	= 1
]
\setmathfont{Libertinus Math}[
	Scale 		= MatchLowercase,
	range		= {scr,bfscr}
]
% Define a custom font face 'mathopfont' and use it to distinguish mathematical operators from prose
\setmathfontface\mathopfont{Latin Modern Roman}[Scale=MatchLowercase]
\setoperatorfont\mathopfont


%::::::::::::::::::::::::::::::::::::::::::::::::::::::::::::::::::::::::::::::::
%                             3. DOCUMENT STRUCTURE                              
%  Article or book-like structure assumed. Here, we just stylise their headers   
%................................................................................
\usepackage{titlesec}		% Controls vertical rhithm and heading aesthetics
\usepackage[titles]{tocloft}	% Change ToC fonts to match chosen heading aesthetics
\usepackage{appendix}		% Manages back-matter structure

% STYLISE HEADINGS
\ifdefined\part
	\titleformat{\part}[display]
		{\sffamily\bfseries\huge\filcenter}	% Format
		{\MakeUppercase{\partname}\ \thepart}	% Label
		{1em}					% Space between number and title
		{\Huge\MakeUppercase}			% Before-code
\fi
\ifdefined\chapter
	\titleformat{\chapter}[display]
		{\sffamily\bfseries\Large}			% Format
		{\MakeUppercase{\chaptertitlename}\ \thechapter}% Label
		{1em}						% Space between number and title
		{\huge\MakeUppercase}				% Before-code
\fi
\titleformat{\section}
    {\sffamily\bfseries\large}	% Format
    {\thesection}		% Label
    {1em}			% Space between number and title
    {\MakeUppercase}		% Before-code
\titleformat{\subsection}
    {\sffamily\bfseries}	% Format
    {\thesubsection}    	% Label
    {1em}               	% Space between number and title
    {\MakeUppercase}    	% Before-code
\titleformat{\subsubsection}
    {\sffamily\bfseries\small}	% Format
    {\thesubsubsection}		% Label
    {1em}               	% Space between number and title
    {\MakeUppercase}    	% Before-code
\titleformat{\paragraph}[runin]
    {\sffamily\bfseries}	% Format
    {\theparagraph}		% Label
    {1em}			% Space between number and title
    {}				% Before-code
\titlespacing*{\paragraph}
    {0pt}			% Left margin indentation
    {3.25ex plus 1ex minus .2ex}% Space BEFORE the paragraph title (vertical)
    {1em}			% Space AFTER the title (horizontal)

%. . . . . . . . . . . . . . . . . . . . . . . . . . . . . . . . . . . . . . . . . 
% 3.2 STYLISE HEADINGS IN ToC SPECIFICALLY
% Set the font for Table of Contents entries to Sans-Serif
\ifdefined\part
	\renewcommand{\cftpartfont}{\sffamily\bfseries}
	\renewcommand{\cftpartpagefont}{\sffamily\bfseries}
\fi
\ifdefined\chapter
	\renewcommand{\cftchapfont}{\sffamily\bfseries}
	\renewcommand{\cftchappagefont}{\sffamily\bfseries}
\fi
\renewcommand{\cfttoctitlefont}{\sffamily\bfseries\large\MakeUppercase}
\renewcommand{\cftsecfont}{\sffamily\bfseries}
\renewcommand{\cftsecpagefont}{\sffamily\bfseries}
\renewcommand{\cftsubsecfont}{\sffamily}
\renewcommand{\cftsubsecpagefont}{\sffamily}
\renewcommand{\cftsubsubsecfont}{\sffamily}
\renewcommand{\cftsubsubsecpagefont}{\sffamily}

%. . . . . . . . . . . . . . . . . . . . . . . . . . . . . . . . . . . . . . . . . 
%\usepackage{fancyhdr}		% Custom headers and footers
%\pagestyle{plain}	% header: empty; footer: centered page number


%::::::::::::::::::::::::::::::::::::::::::::::::::::::::::::::::::::::::::::::::
%                                    4. FLOATS                                   
%               Manage Floating Content (e.g.: figures and tables)               
%................................................................................
% FLOATS IN GENERAL
\usepackage{float}	% Provides [H] 'Here' float placement specifier (use with caution)
\usepackage{subcaption}	% Format multiple images (tables) in one figure (table) environment
			% NOTE: This package *automatically* loads the 'caption' package,
\captionsetup{
	labelfont = {small,bf},
	textfont  = {small,it}
}
%\subcaptionsetup{...}	% Uncomment to customize sub-captions separately

%. . . . . . . . . . . . . . . . . . . . . . . . . . . . . . . . . . . . . . . . . 
% TABLES
\usepackage{booktabs}	% Format tables following scientific convention
\usepackage{multirow}	% Allow multirow cells in tables

%. . . . . . . . . . . . . . . . . . . . . . . . . . . . . . . . . . . . . . . . . 
% IMAGES
% NATIVE PROGRAMMATIC DIAGRAMS & PLOTS
\usepackage{tikz}		% Create vector graphics programmatically
\usetikzlibrary{arrows.meta}	% Additional arrow styles
\usetikzlibrary{calc}		% Coordinate calculations
\usetikzlibrary{positioning}	% Better relative placement of notes
\usepackage{pgfplots}		% Tool for creating function plots and charts
\pgfplotsset{compat=1.18}	% Enforce version 1.18 for consistency

% EXTERNAL IMAGE MANAGEMENT
\usepackage{graphicx}		% LaTeX standard for including external image files

%. . . . . . . . . . . . . . . . . . . . . . . . . . . . . . . . . . . . . . . . . 
% ALGORITHMS
\usepackage[
	portuguese,
	ruled,
	vlined,
	linesnumbered
]{algorithm2e}	% Typeset formal algorithms
% Set algorithm font to the chosen monospace font
\SetAlFnt{\small\ttfamily}
\SetAlCapFnt{\small\ttfamily}
% Define Portuguese keywords
\SetKwInput{KwRequer}{Requer}
\SetKwInput{KwPrecondicao}{Pré-condição}
\SetKwInput{KwInicializa}{Inicializa}
\SetKwFor{Para}{para}{faça}{fim para}
\SetKwIF{Se}{SenaoSe}{Senao}{se}{então}{senão se}{senão}{fim se}
\DontPrintSemicolon


%::::::::::::::::::::::::::::::::::::::::::::::::::::::::::::::::::::::::::::::::
%                             5. MICRO-LAYOUT                              
%                       Packages that alter "inline" text                        
%................................................................................
% TEXT
\usepackage{multicol}		% Environment with multiple columns of text
\usepackage{enumitem}		% Customize list environments (provides 'resume')

%. . . . . . . . . . . . . . . . . . . . . . . . . . . . . . . . . . . . . . . . . 
% Now... for the micro-layout good stuff...
%
% █▀▄▀█ ▄▀█ ▀█▀ █ █
% █ ▀ █ █▀█  █  █▀█
%
% NICHE MATH PACKAGES
\usepackage{xfrac}	% Slanted fractions (\sfrac)
\usepackage{cancel}	% Strike-through for symbols (\cancel and \cancelto)
\usepackage{amsthm}	% AMS's package for creating theorem-like environments

%::::::::::::::::::::::::::::::::::::::::::::::::::::::::::::::::::::::::::::::::::
%                         5.1 CUSTOM THEOREM ENVIRONMENTS                          
%                                   (Colour-coded)                                  
%..................................................................................
% DEFINE COLOUR CODE FOR EACH THEOREM-LIKE ENVIRONMENT;
% (makes changing the colour choice effortless!)
\colorlet{lmm}{DarkOrchid}	% For lemma env.
\colorlet{thm}{Tan}		% For theorem env.
\colorlet{prop}{Maroon}		% For proposition env.
\colorlet{crl}{RawSienna}	% For corollary env.
\colorlet{def}{OliveGreen}	% For definition env.
\colorlet{eg}{RoyalPurple}	% For example env.
\colorlet{ex}{RoyalPurple}	% For exercise env.
\colorlet{obs}{MidnightBlue}	% For remark env.
\colorlet{note}{PineGreen}	% For notation remark env.

%. . . . . . . . . . . . . . . . . . . . . . . . . . . . . . . . . . . . . . . . . 
% DEFINE CUSTOM THEOREM STYLES
% > STYLE I - TAUTOLOGIES - FIRST NUMBER, THEN NAME
\newtheoremstyle{number-name}
{4mm}						% Space above env.
{4mm}						% Space below env.
{\itshape}					% Body font
{}						% Indent amount
{\bfseries}					% Head font
{:}						% Punctuation after head
{.3em}						% Space after head
{\thmnumber{#2}\,\thmname{#1}\thmnote{ (#3)}}	% Head spec

% > STYLE II - EXAMPLES & EXERCISES - FIRST NAME, THEN ITALIC NUMBER 
\newtheoremstyle{name-it_number}
{4mm}						% Space above env.
{4mm}						% Space below env.
{\itshape}					% Body font
{}						% Indent amount
{\bfseries}					% Head font
{}						% Punctuation after head
{.3em}						% Space after head
{\thmname{#1}\,{\itshape\thmnumber{#2}}}	% Head spec

% > STYLE III - REMARKS, SOLUTIONS & PROOFS - NO NUMBERING (should refer to a numbered env.)
\newtheoremstyle{no_number}
{4mm}						% Space above env.
{4mm}						% Space below env.
{}						% Body font
{}						% Indent amount
{\bfseries\itshape}				% Head font
{.}						% Punctuation after head
{.3em}						% Space after head
{}						% Head spec

% > STYLE IV - NOTATION REMARKS - NO NUMBERING (should refer to a numbered env.) AND ITALIC BODY 
\newtheoremstyle{no_number-it_body}
{4mm}						% Space above
{4mm}						% Space below
{\itshape}					% Body font
{}						% Indent amount
{\bfseries}					% Head font
{:}						% Punctuation after head
{.3em}						% Space after head
{}						% Head spec

%. . . . . . . . . . . . . . . . . . . . . . . . . . . . . . . . . . . . . . . . .
% APPLY THEOREM STYLES AND DEFINE RESPECTIVE ENVIRONMENTS
% > STYLE I - TAUTOLOGIES - FIRST NUMBER, THEN NAME
\theoremstyle{number-name}
  \newtheorem{lemma}{\color{lmm}Lema}[section]
  \newtheorem{theorem}{\color{thm}Teorema}[section]
  \newtheorem{proposition}{\color{prop}Proposição}[section]
  \newtheorem{corollary}{\color{crl}Corolário}[section]
  \newtheorem{definition}{\color{def}Definição}[section]

% Manually force counters to share the 'lemma' counter (i.e.: shared counter for style I)
\makeatletter
	\let\c@theorem\c@lemma
	\let\c@proposition\c@lemma
	\let\c@corollary\c@lemma
	\let\c@definition\c@lemma
\makeatother

% > STYLE II - EXAMPLES & EXERCISES - FIRST NAME, THEN ITALIC NUMBER 
\theoremstyle{name-it_number}
  \newtheorem{exercise}{\textcolor{ex}{Exercício}}[section]
  \newtheorem{example}{\textcolor{eg}{Exemplo}}[section]

% Manually force 'example' to use the 'exercise' counter (i.e. shared counter for style II)
\makeatletter
\let\c@example\c@exercise
\makeatother

% > STYLE III - REMARKS, SOLUTIONS & PROOFS - NO NUMBERING
\theoremstyle{no_number}
  % Remark
  \newtheorem*{observation*}{\color{obs}Observação}
  % Solution
  \newtheorem*{solution}{\color{ex}Solução}
  % Proofs
  \newtheorem*{proofT}{\color{thm}Demonstração}
  \newtheorem*{proofL}{\color{lmm}Demonstração}
  \newtheorem*{proofC}{\color{crl}Demonstração}
  \newtheorem*{proofP}{\color{prop}Demonstração}
  \newtheorem*{proofE}{\color{eg}Demonstração}
  \newtheorem*{proofO}{\color{obs}Demonstração}

% > STYLE IV - NOTATION REMARKS - NO NUMBERING AND ITALIC BODY 
\theoremstyle{no_number-it_body}
  \newtheorem*{notation*}{\color{note}Notação}

%. . . . . . . . . . . . . . . . . . . . . . . . . . . . . . . . . . . . . . . . .
% EXTRA: CUSTOM QED SYMBOL (a small 'black square' U+25A0)
\renewcommand\qedsymbol{■}


%:::::::::::::::::::::::::::::::::::::::::::::::::::::::::::::::::::::::::::::::::
%                            5.2 CUSTOM MATH COMMANDS                             
%.................................................................................
\newcommand{\Bvec}[1]{\symbf{#1}}	% Bold symbol vector notation

% CONSTANTS
\newcommand{\e}{\symup{e}}	% ISO notation for Euler's number
\newcommand{\im}{\symup{i}}	% ISO notation for the imaginary unit
\newcommand{\piup}{\symup{\pi}}	% ISO notation for pi

% RELATIONS
\newcommand{\longto}{\longrightarrow}	% Alias to note function domain/codomain (f: A --> B)

% Below, a colon-equals (and an analogous equals-colon) sign built 'from scratch', to be more compact and elegant
\newcommand*{\Defeq}{%
	\mathrel{%				% Wrapper so TeX treats/interprets the symbol as a relation
		\vcenter{%			% Vertically center the 'colon' part
			\baselineskip0.5ex%	% Manually set the vertical spacing between the dots
			\lineskiplimit0pt%	% Disable TeX's 'anti-collision' rule for full control
			\hbox{\scriptsize.}%	% The top dot
			\hbox{\scriptsize.}%	% The bottom dot
		}%
		{=}%				% The equals sign
	}%
}					% Compact colon-equals
\newcommand*{\Eqdef}{%
	\mathrel{%				% Wrapper so TeX treats/interprets the symbol as a relation
		{=}%				% The equals sign
		\vcenter{%			% Vertically center the 'colon' part
			\baselineskip0.5ex%	% Manually set the vertical spacing between the dots
			\lineskiplimit0pt%	% Disable TeX's 'anti-collision' rule for full control
			\hbox{\scriptsize.}%	% The top dot
			\hbox{\scriptsize.}%	% The bottom dot
		}%
	}%
}					% Compact equals-colon
\newcommand{\subnormal}{\triangleleft}	% Normal subgroup (left)
\newcommand{\supnormal}{\triangleright}	% Normal subgroup (right)

% UNARY OPERATIONS
\newcommand{\T}{{\symsfup{T}}}	% Transpose symbol
\newcommand{\Tp}[1]{{#1}^{\T}}	% Transpose operation

% OPERATORS
\newcommand{\diff}{\mathop{}\!\symup{d}}% Differential operator (with spacing)
\DeclareMathOperator{\dist}{dist}
\DeclareMathOperator{\map}{map}

% LOGIC
\newcommand{\limplies}{\rightarrow}	% Conditional connector
\newcommand{\liff}{\leftrightarrow}	% Biconditional connector
\newcommand{\nand}{\uparrow}		% NAND connector
\newcommand{\nor}{\downarrow}		% NOR connector
\newcommand{\xor}{\veebar}		% XOR connector


%:::::::::::::::::::::::::::::::::::::::::::::::::::::::::::::::::::::::::::::::::
%                                  FINALIZATION                                   
%                 Packages that patch others. MUST BE LOADED LAST                 
%.................................................................................
% BIBLIOGRAPHY                                   
\usepackage[
	backend	= biber,	% Use Biber backend (more powerful than BibTex)
	style	= ieee,
	backref	= true
]{biblatex}

%. . . . . . . . . . . . . . . . . . . . . . . . . . . . . . . . . . . . . . . . .
% HYPERLINKS
\usepackage{hyperref}	% Add hyperlinks to references, ToC and URLs
\hypersetup{
	plainpages	= false,	% Attempts to prevent duplicate page anchors
	colorlinks	= true,		% true: colored links; false: boxed links
	linkcolor	= NavyBlue,	% Internal links
	citecolor	= NavyBlue,
	urlcolor	= NavyBlue,
	breaklinks	= true,
	pdfborder	= {0 0 0}
}

%. . . . . . . . . . . . . . . . . . . . . . . . . . . . . . . . . . . . . . . . .
% REFERENCING
\usepackage[
	brazilian,
	nameinlink
]{cleveref}	% \cref is objectively better than \autoref

% DEFINE 'cleveref' NAMES FOR CUSTOM THEOREM-LIKE ENVIRONMENTS
	% STYLE I ENVIRONMENTS
\crefname{lemma}{lema}{lemas}
\Crefname{lemma}{Lema}{Lemas}
\crefname{theorem}{teorema}{teoremas}
\Crefname{theorem}{Teorema}{Teoremas}
\crefname{proposition}{proposição}{proposições}
\Crefname{proposition}{Proposição}{Proposições}
\crefname{corollary}{corolário}{corolários}
\Crefname{corollary}{Corolário}{Corolários}
\crefname{definition}{definição}{definições}
\Crefname{definition}{Definição}{Definições}
	% STYLE II ENVIRONMENTS
\crefname{exercise}{exercício}{exercícios}
\Crefname{exercise}{Exercício}{Exercícios}
\crefname{example}{exemplo}{exemplos}
\Crefname{example}{Exemplo}{Exemplos}

\newcommand{\customref}[2]{%	% Simple alias for custom link text
	\hyperref[#2]{#1}%
}

%. . . . . . . . . . . . . . . . . . . . . . . . . . . . . . . . . . . . . . . . .
% MICRO-TYPOGRAPHY
\usepackage{microtype}		% Enhance micro-typography (e.g.: protrusion, expansion)
				% Loaded last for best compatibility with other packages
